\chapter*{Abstract}
\markboth{\small \textit{Abstract}}{\small \textit{Abstract}}
\addcontentsline{toc}{chapter}{Abstract}

Il modello Function-as-a-Service (FaaS) ha acquisito crescente rilevanza come
paradigma di esecuzione event-driven nel continuum
Edge--Cloud: distribuire funzioni
stateless su nodi eterogenei e geograficamente distribuiti consente di ridurre la
latenza percepita e soddisfare i requisiti di applicazioni real-time.
Tuttavia, le piattaforme FaaS tipicamente guidano le decisioni di scheduling in
base a latenza, disponibilità e uso delle risorse, senza considerare in modo
esplicito il costo carbonico dell'energia consumata.
In un sistema geograficamente distribuito, la carbon intensity della rete
elettrica, ovvero la quantità di emissioni di CO\textsubscript{2} associata
alla produzione o al consumo di una determinata quantità di energia elettrica,
varia nel tempo e nello spazio in funzione del mix energetico locale; di
conseguenza, una stessa invocazione può produrre emissioni molto diverse a
seconda del luogo e del momento in cui viene eseguita. Il problema affrontato
in questa tesi è quindi come ridurre l'impatto carbonico operativo delle
funzioni FaaS senza imporre una degradazione indiscriminata della qualità del
risultato.

Allo stato dell'arte, il problema della sostenibilità nei sistemi serverless
viene affrontato lungo direzioni complementari ma parziali. Le tecniche di
approximate computing sfruttano l'esistenza di implementazioni alternative di
una stessa funzione, ciascuna con un diverso profilo energia--qualità, per
modulare tale compromesso a seconda delle condizioni di esecuzione, ma senza
collegarlo alla variabilità del segnale carbonico. Le politiche carbon-aware
considerano le emissioni di carbonio agendo sul posizionamento geografico del
carico o sul suo differimento temporale. Rimane quindi aperta la possibilità di
selezionare, a livello locale, quale implementazione di una funzione eseguire
in funzione dell'impatto carbonico corrente e del compromesso tra energia e
qualità.

Il risultato principale di questa tesi è un'estensione carbon-aware della
piattaforma open-source Serverledge, nella quale ogni funzione logica è
associata a un insieme di varianti funzionalmente equivalenti ma differenziate
per profilo energia--qualità. Uno scheduler adattivo seleziona a runtime la
variante più adatta alla quantità corrente di emissioni di CO\textsubscript{2},
privilegiando l'accuratezza nelle fasce di bassa intensità e orientandosi
verso varianti a minore consumo quando la rete è più carbon-intensive. Il
sistema realizza così un compromesso adattivo, evitando sia una politica che
fornisce sempre la versione più accurata sia una politica sempre orientata al
minimo consumo.

La metodologia si articola in tre componenti integrati. I profili energia--qualità
delle varianti vengono stimati a runtime tramite Kepler, strumento di misura del
consumo energetico a granularità di container, e Prometheus, sistema di raccolta
di metriche; tali profili vengono aggiornati incrementalmente tramite media mobile
esponenziale per adattarsi al comportamento osservato sul nodo. Poiché varianti
poco invocate possono avere profili stimati con scarsa precisione, la selezione
incorpora un termine di esplorazione ispirato al Lower Confidence Bound, che evita
di penalizzarle prematuramente. Sulla base dei profili correnti, lo scheduler
costruisce la frontiera di Pareto, limitando la scelta alle varianti non
dominate, cioè a quelle per cui non esiste un'alternativa simultaneamente
migliore in termini di consumo e qualità. Il
segnale carbonico, acquisito tramite l'API ElectricityMaps e normalizzato rispetto
ai percentili storici della zona attraverso la \emph{Carbon Response Function}
(CRF), una funzione logistica che restituisce un parametro $\lambda \in [0,1]$
indicante quanto la carbon intensity corrente si discosti dalla media storica
locale, viene usato per scalarizzare tale frontiera: valori elevati di $\lambda$
orientano la selezione verso varianti più accurate, mentre valori bassi
favoriscono varianti a minore consumo energetico.

Il contributo distintivo di questo lavoro consiste nel combinare profilazione
energetica adattiva a livello di container, selezione di varianti approssimate
orientata alla qualità e un segnale carbonico calibrato per zona in una
piattaforma FaaS edge-oriented. A differenza delle soluzioni esistenti, la
politica proposta non si limita a osservare i consumi né ad agire sul
posizionamento del carico, ma costruisce a ogni invocazione un compromesso
dipendente dall'impatto carbonico corrente e dai profili osservati sul nodo,
senza ricorrere a politiche statiche.

La valutazione sperimentale su cinque zone
europee e sei funzioni, tre numeriche e tre di classificazione ML, mostra che,
rispetto alla politica che seleziona sempre la variante a massima qualità, il
sistema ottiene risparmi energetici tra il 62\% e il 92\% e una riduzione
stimata delle emissioni operative di CO\textsubscript{2} fino al 94{,}82\%,
con un degrado qualitativo contenuto tra lo 0{,}01\% e lo 0{,}61\% per le
funzioni numeriche e tra il 4{,}34\% e il 10{,}55\% per le funzioni ML.
Rispetto alla baseline che privilegia sempre il minimo consumo, lo scheduler
garantisce un vantaggio qualitativo significativo per le funzioni ML,
con guadagni tra il 23{,}40\% e il 32{,}19\%. Prove con carico crescente confermano che il sistema rimane stabile
sotto pressione: in presenza di fenomeni di burst, il maggiore costo
computazionale delle varianti ad alta qualità può accelerare la saturazione
del server, mentre l'approccio proposto riduce la pressione computazionale
media sul nodo quando la quantità corrente di emissioni di CO\textsubscript{2}
rende preferibili varianti meno costose.

I risultati confermano che è possibile ridurre significativamente l'impatto
carbonico operativo delle funzioni FaaS integrando la consapevolezza delle
emissioni direttamente nella politica di scheduling, senza imporre una
degradazione indiscriminata della qualità del servizio.
